\documentclass[conference]{IEEEtran}
\IEEEoverridecommandlockouts

\usepackage[utf8]{inputenc}
\usepackage{amsmath,amssymb,amsfonts}
\usepackage{algorithmic}
\usepackage{algorithm}
\usepackage{graphicx}
\usepackage{booktabs}
\usepackage{multirow}
\usepackage[caption=false,font=footnotesize]{subfig}
\usepackage{xcolor}
\usepackage{cite}
\usepackage{url}
\usepackage{bm}
\usepackage{float}
\usepackage{tikz}
\usetikzlibrary{positioning,arrows.meta,fit,backgrounds}

\graphicspath{{figures/}}
\raggedbottom
% ============================================================
\title{Digital Twins and Multi-Agent Reinforcement Learning for IoT Resource Allocation}
% ============================================================
\author{
Thi Huyen Trang Nguyen\textsuperscript{1}, 
Trong Binh Hoang\textsuperscript{1}, 
Quoc Khanh Nguyen\textsuperscript{1},\\ 
Van Khuong Ngo\textsuperscript{2}, 
Van Tiep Pham\textsuperscript{2}, and 
Van Linh Dinh\textsuperscript{2} \\[1em]

\normalsize
\textsuperscript{1}Le Quy Don Technical University, Hanoi, Vietnam\\
\textsuperscript{2}Academy of Cryptography Techniques, Hanoi, Vietnam\\[0.5em]

\textsuperscript{*}Corresponding authors: vanlinh@actvn.edu.vn, binhht99@lqdtu.edu.vn, trangcucheng@gmail.com}
% ============================================================
\begin{document}
\maketitle
\begin{abstract}
% ============================================================
Effective IoT network management requires simultaneous semantic feature
compression, real-time channel quality monitoring, anomaly detection, and
adaptive resource allocation. Existing approaches address these tasks in
isolation, lacking a unified system that jointly exploits temporal dynamics
and multi-agent coordination. In this paper, a three-stage framework is proposed that integrates
Variational Autoencoder (VAE)-based semantic compression, an Extended Kalman Filter (EKF)-based Digital Twin for joint
quality monitoring and anomaly detection, and Multi-Agent Deep Reinforcement Learning (MADRL) with Graph Attention
Networks (MADRL-GAT) for cooperative resource allocation.
Experiments on three heterogeneous real-world IoT datasets
including Sigfox Antwerp, LoRaWAN Italy, and Indoor demonstrate that the proposed
EKF-based Digital Twin substantially outperforms unsupervised clustering baselines
(K-Means, GMM) in quality classification and delivers competitive anomaly detection
against dedicated single-task detectors including LOF, Isolation Forest, Z-Score, while
being the only evaluated method to jointly deliver both capabilities in a single
recursive filter pass without requiring labeled fault data. MADRL-GAT achieves
stable policy convergence across all settings, and the full pipeline operates in
real time on CPU, confirming practical suitability and cross-technology
generalization across heterogeneous IoT deployments.
\end{abstract}

\begin{IEEEkeywords}
Internet of Things, digital twin, multi-agent reinforcement learning, IoT resource allocation, IoT security
\end{IEEEkeywords}

\section{Introduction}
\label{sec:intro}

The explosive growth of IoT deployments, which spans smart
cities, industrial sensing, and indoor positioning, has introduced
unprecedented demands on wireless network infrastructure~\cite{ref_iot_overview}. Devices operating
across heterogeneous technologies such as Sigfox, LoRaWAN, and WiFi
continuously generate high-dimensional telemetry streams that must be
processed, monitored, and acted upon with minimal latency and energy
overhead. Two fundamental challenges arise: the raw
feature space is high-dimensional and noisy, making downstream learning
inefficient, and the network state evolves continuously over time, subject to both
signal quality degradation and potential adversarial anomalies. Without
mechanisms that simultaneously monitor connection quality and detect
abnormal behaviour, resource allocation decisions risk being
systematically misguided or exploited.

A natural approach to the dimensionality challenge is to extract compact,
task-relevant representations from raw sensor observations before any
downstream processing. Liu et al.~\cite{ref_vae_wireless} proposed a
task-adaptive semantic compression scheme employing a VAE to encode raw sensor observations into compact representations,
substantially reducing communication overhead in resource-constrained
networks. The compression rate adapts to instantaneous channel quality,
yielding strong reconstruction fidelity across diverse wireless conditions.
However, this scheme targets static point-to-point links without temporal
state tracking, provides no anomaly-detection capability, and does not
coordinate multiple agents, which restricts its applicability to the dynamic,
heterogeneous IoT settings addressed in this work.

Beyond feature compression, Deep Reinforcement Learning (DRL) has emerged as a leading approach to
adaptive resource management under time-varying interference, with
multi-agent formulations enabling distributed control across heterogeneous
wireless settings~\cite{ref_drl_comm2}. In the IoT domain,
Xiong et al.~\cite{ref_drl_iot} proposed a DRL-based resource allocation framework for IoT edge computing,
demonstrating adaptive multi-resource management under dynamic task arrivals and heterogeneous device requirements.
A key limitation of this approach is that
raw, high-dimensional feature vectors serve as the agent state without
semantic compression, inflating the effective state space and degrading
sample efficiency. Moreover, no channel-quality feedback enters the reward
signal, preventing the agent from distinguishing nominal from degraded
link conditions. These shortcomings motivate coupling feature compression
with a predictive channel model before resource allocation.
More recently, Wali et al.~\cite{ref_wali2025} proposed a semantic-aware hybrid framework on the same
Sigfox and LoRaWAN datasets, integrating a single-agent DQN with supervised Random Forest
classification and Isolation Forest anomaly detection.
While this work demonstrates the value of combining semantic features with ensemble learning,
it relies on labeled signal-quality annotations for the classification stage, operates with a
single RL agent without inter-agent coordination, and does not employ a principled temporal
state estimator for channel monitoring. These constraints limit its applicability in
label-scarce, multi-device IoT deployments, motivating the fully unsupervised and
multi-agent design proposed in this work.

Digital twin (DT) technology has emerged as a powerful paradigm for
real-time monitoring and prediction in IoT deployments, enabling
softwarized replicas of physical systems that continuously synchronize
with their physical counterparts~\cite{ref_dt_iot}.
Applied to wireless resource management, DT-based approaches use predicted
channel states to guide allocation decisions ahead of reactive baselines.
Among these, Tang et al.~\cite{ref_dt_resource} combined a DT with distributed deep reinforcement learning
for resource allocation in industrial IoT network slicing, achieving predictive resource control that
substantially outperforms non-predictive baselines. Nevertheless,
quality classification in such systems relies on fixed decision thresholds that cannot
adapt to heterogeneous dataset statistics, anomaly detection is absent,
and raw features are fed directly to the allocation module without
semantic compression, leaving these systems unable to flag silent sensor
faults or generalize across heterogeneous channel conditions.

The most closely related integrated framework is that of Hazarika
et al.~\cite{ref_semadt}, which combined semantic compression, predictive
DT forecasting, and MADRL with a graph neural network (GNN)
for maritime search-and-rescue IoT operations, demonstrating that semantic
state inputs and DT-predicted quality signals together substantially outperform
non-predictive Multi-Agent Reinforcement Learning (MARL) baselines. Despite this advance, the Graph Convolutional Network (GCN) layer applies
symmetric neighborhood aggregation without capturing heterogeneous inter-agent
relationship strengths, whereas Graph Attention Networks (GAT)~\cite{ref_gat_comm}
do so via learned attention weights, enabling each agent to selectively
emphasize the most relevant neighbors (e.g., strongly interfering or
heavily loaded nodes) and thereby make more informed resource allocation
decisions in heterogeneous IoT deployments. In addition, the anomaly detector uses a single hand-picked scalar
cutoff applied to a univariate signal, which ignores the joint
covariance structure of the latent state vector and lacks the
statistical guarantees of a multivariate test such as the
chi-squared-based Normalized Innovation Squared (NIS) criterion~\cite{ref_ekf}.
Finally, experiments are conducted exclusively on a maritime
search-and-rescue scenario, so it remains unknown whether the
approach generalises to other IoT technologies such as LoRaWAN
or WiFi-based sensor networks.

In general, three critical gaps remain in the existing literature. First,
the broader literature largely treats semantic compression, DT
forecasting, and multi-agent resource allocation as independent
modules~\cite{ref_vae_wireless,ref_dt_iot,ref_dt_resource}. Even the
most integrated recent system~\cite{ref_semadt} lacks an explicit feedback
path in which anomaly signals from the DT are fed back into
the resource allocation reward. Second, existing DT implementations
address either quality classification or anomaly detection in
isolation~\cite{ref_dt_resource,ref_anomaly_survey}, rather than jointly performing
both within a single principled state estimator. Third, cross-technology
generalizability across heterogeneous IoT protocols (Sigfox, LoRaWAN, WiFi)
remains unexplored, as prior integrated systems have been validated only
in a single domain~\cite{ref_semadt}.

In this paper, a unified three-stage intelligent framework is proposed to
simultaneously address the above gaps. To close the first gap, a
VAE establishes a continuous, differentiable
feedback path from raw IoT measurements to compact semantic latent codes
that are shared across all downstream stages. To close the second gap,
an EKF-based DT jointly performs
connection quality classification and anomaly detection within a single
principled state estimator, coupling both outputs into a unified signal.
To close the third gap, a MADRL-GAT module coordinates
resource allocation across heterogeneous IoT devices by consuming the
semantic codes, quality labels, and anomaly flags produced upstream,
and is validated across three distinct wireless technologies. The proposed framework is evaluated on three real-world IoT datasets
spanning heterogeneous wireless technologies: Sigfox
Antwerp~\cite{ref_sigfox}, LoRaWAN Italy~\cite{ref_lorawan}, and
Indoor WiFi~\cite{ref_indoor}. Results demonstrate that the EKF-based
DT consistently outperforms classical clustering baselines in
both quality classification and anomaly detection, while the MADRL-GAT
module converges stably and achieves high cumulative rewards across all
three datasets, confirming the effectiveness of the integrated
semantic-and quality-aware design.

The main contributions are summarized as follows:
\begin{itemize}
\item A unified three-stage pipeline is proposed that couples VAE semantic
compression, EKF DT monitoring, and MADRL-GAT resource allocation,
with compressed latent codes serving as both agent state and quality-aware
reward input.
\item An EKF-based DT is designed to jointly perform quality
classification and anomaly detection in a single recursive estimation pass.
Quality thresholds are set automatically via Youden's J~\cite{ref_youden}, and anomalies are
scored via the NIS chi-squared test, providing statistically principled fault
detection absent in prior fixed-threshold DT systems.
\item Experiments on three heterogeneous real-world IoT datasets (Sigfox
Antwerp, LoRaWAN Italy, Indoor) show that the EKF DT outperforms
K-Means and GMM on quality classification AUC across all datasets, and
achieves anomaly detection AUC of 0.987 on Indoor, while MADRL-GAT converges
stably within 50 episodes.
\end{itemize}

The remainder of this paper is organised as follows. Section~\ref{sec:model}
presents the system model. Section~\ref{sec:problem} formulates the
optimization problem. Section~\ref{sec:method} details the proposed framework.
Section~\ref{sec:experiments} reports experimental results.
Section~\ref{sec:conclusion} concludes the paper.

\section{System Model}
\label{sec:model}
% ============================================================

\subsection{IoT Network Model}

\begin{figure}[t]
  \centering
  \includegraphics[width=\columnwidth]{figures/iotnodes.pdf}
  \caption{Heterogeneous IoT network model with three wireless technologies.}
  \label{fig:iotnodes}
\end{figure}

The system consists of a heterogeneous IoT network comprising $N$ nodes
indexed by $i\in\{1,\ldots,N\}$, where each node periodically transmits
sensor readings over a wireless channel to a gateway or base station. Let
$\mathbf{x}_i(t)\in\mathbb{R}^d$ denote the raw feature vector
observed by node $i$ at time $t$, comprising signal strength measurements
and sensor readings (e.g., RSSI per base station, SNR, temperature,
humidity, battery level).

\subsection{Channel Quality State Model}

The quality of the wireless link for node $i$ is characterized by two
key indicators: RSSI and the number of active base stations $n^{\rm BS}_i(t)$.
These form the two-dimensional channel state
$\mathbf{s}_i(t)=[{\rm RSSI}_i(t),\;n^{\rm BS}_i(t)]^\top$, which is
the input to the DT-EKF (Section~\ref{sec:dt}).
Since no ground-truth quality annotations are available, pseudo binary
labels are constructed from the raw RSSI for evaluation purposes only.
A sample is labeled good quality if ${\rm RSSI}_i(t)\geq\tau_i$.
The natural class distribution is severely imbalanced: for Sigfox Antwerp,
only 1.9\% of samples exceed the standard Sigfox link-budget threshold of
$-$130~dBm, rendering standard accuracy and Area Under the Curve (AUC) metrics unreliable.
To obtain stable, balanced evaluation labels, $\tau_i$ is therefore set
to the sample median of RSSI observations in each dataset, which
guarantees approximately 50\% positive labels by construction regardless
of the underlying distribution shape. These median-derived pseudo-labels
serve exclusively as the reference for computing AUC and accuracy in
Table~\ref{tab:main_results}. Separately, the DT-EKF derives
its own classification threshold $\tau^*$ via Youden's J
(Section~\ref{sec:dt}), optimizing True Positive Rate (TPR)$-$False Positive Rate (FPR) on the filtered state
estimate $\widehat{\rm RSSI}_{t|t}$ against these same pseudo-labels.

\subsection{Resource Allocation and Communication Model}

At each decision epoch, node $i$ selects a discrete action
$a_i(t)\in\{1,\ldots,A\}$ representing its resource configuration
(e.g., transmission power level, spreading factor, or channel index).
Because the three real-world datasets provide empirical RSSI measurements
rather than explicit channel coefficients, the normalized raw RSSI
$\bar{s}_i(t)\in[0,1]$ serves as a tractable, scale-invariant proxy
for link quality in the reward signal, as detailed in
Section~\ref{sec:method}. Nodes are connected through a graph
$\mathcal{G}=(\mathcal{V},\mathcal{E},\mathbf{A})$, where the adjacency
matrix $\mathbf{A}$ encodes inter-node relationships.
For Sigfox Antwerp, $\mathbf{A}$ is derived from RSSI cross-correlation
among the top-10 most active base stations, reflecting co-coverage
structure where correlated nodes likely share overlapping BS footprints
and mutually cause interference.
For LoRaWAN and Indoor, $\mathbf{A}$ is set to fully-connected, as the
low input dimensionality (5 features) provides insufficient basis for
reliable pairwise correlation estimation, and the small agent count
makes full connectivity computationally negligible.
Fig.~\ref{fig:graph_adj} illustrates the two adjacency construction policies.

\begin{figure}[t]
  \centering
  % Compile fig_graph_construction.tex as standalone PDF,
  % then upload fig_graph_construction.pdf to the figures/ folder on Overleaf.
  \includegraphics[width=\columnwidth]{fig_graph_construction}
  \caption{Graph adjacency construction for
  $\mathcal{G}=(\mathcal{V},\mathcal{E},\mathbf{A})$.}
  \label{fig:graph_adj}
\end{figure}

\subsection{Fault Model}

To evaluate anomaly detection in the absence of naturally labeled faults,
synthetic anomalies are injected into 5\% of each time-ordered sequence
via three fault types, each shifting the RSSI within $\pm2\sigma$ of the
global distribution to avoid trivially detectable outliers.
The 5\% injection rate follows standard anomaly detection practice~\cite{ref_anomaly_survey}
and reflects realistic IoT fault prevalence, where silent hardware and link faults
typically account for 1--5\% of transmissions in deployed networks.
Fault magnitudes are bounded at $\pm2\sigma$ because this boundary corresponds
to the 95th percentile of a Gaussian distribution, placing injected values
at the edge of the plausible data envelope rather than in the extreme tail.
Exceeding this bound (e.g., $\pm3\sigma$ or beyond) would produce
outliers detectable by any simple threshold rule, eliminating the ability
to differentiate methods. A reduced magnitude of $1.5\sigma$ is used for
the stuck fault, which represents gradual sensor drift rather than an
abrupt event. The three fault types are:
\begin{itemize}
  \item Spike (interference): RSSI is raised to
    $\mu_{\rm RSSI}+2\sigma_{\rm RSSI}$, simulating a sudden surge
    caused by co-channel interference.
  \item Drop (attenuation): RSSI is lowered to
    $\mu_{\rm RSSI}-2\sigma_{\rm RSSI}$ and $n^{\rm BS}_i$ is set to
    zero, simulating abrupt signal loss.
  \item Stuck (sensor drift): RSSI is fixed at
    $\mu_{\rm RSSI}+1.5\sigma_{\rm RSSI}$ and $n^{\rm BS}_i$ is set to
    zero, simulating a sensor that stops updating its reading.
\end{itemize}

% ============================================================
\section{Problem Formulation}
\label{sec:problem}
% ============================================================

Solving simultaneously for the VAE parameters $(\phi,\theta)$
and the multi-agent actions $\{a_i\}$ in a single joint pass is
computationally intractable and prone to gradient instability. The
problem is therefore decomposed into a \textit{sequential two-phase
paradigm} (Fig.~\ref{fig:framework}).
Phase~1 trains the VAE offline to maximize the standard Evidence Lower BOund (ELBO)~\cite{ref_vae},
then freezes the encoder. Phase~2 optimizes the per-agent actions
$\{a_i(t)\}$ via the decentralized RL problem:
\begin{equation}
(\mathcal{P}_{\rm RL})\!:\;\max_{\{a_i\}}\; \mathbb{E}\!\left[\sum_{t=0}^{T}\sum_{i=1}^{N} r_i(t) \;\Bigg|\; \hat{\mathbf{z}}_i(t)=\bm{\mu}_{\phi^*}(\mathbf{x}_i(t))\right]
\label{eq:rl_opt}
\end{equation}
where $r_i(t)$ is the per-agent reward defined in Section~\ref{sec:method},
quality labels satisfy $\hat{q}_i(t)=\mathbf{1}[\widehat{\rm RSSI}_{i,t}\geq\tau^*]$
with $\tau^*$ from Youden's J (Section~\ref{sec:dt}), and
$a_i(t)\in\{1,\ldots,A\}$. This decoupled formulation guarantees
convergence and practical feasibility.

% ============================================================
\section{Proposed Framework}
\label{sec:method}
% ============================================================

\begin{figure*}[t]
  \centering
  \includegraphics[width=\textwidth]{figures/lưu_đồ_thuật_toán (3).pdf}
  \caption{Overview of the proposed framework.}
  \label{fig:framework}
\end{figure*}

Fig.~\ref{fig:framework} illustrates the overall architecture. The VAE
is pre-trained offline. At runtime, the frozen encoder projects each
observation $\mathbf{x}_i(t)$ to a semantic latent code
$\hat{\mathbf{z}}_i(t)$ that feeds the MADRL agents. Concurrently, the
observation sequence $\mathbf{s}_i(t)$ drives the DT-EKF, whose
quality labels shape the agent reward signal. The three stages are detailed
below.

\subsection{Stage 1: Semantic Feature Compression via VAE}

The VAE is trained offline following the standard ELBO
objective~\cite{ref_vae}. At inference, the frozen encoder maps each
raw observation $\mathbf{x}_i(t)\in\mathbb{R}^d$ to a compact
deterministic latent code
$\hat{\mathbf{z}}_i(t)=\bm{\mu}_\phi(\mathbf{x}_i(t))\in\mathbb{R}^{d_z}$,
which serves as the MADRL agent state. The latent dimension $d_z$ is
the smallest value that keeps reconstruction mean squared error (MSE) within an acceptable
bound: $84\to16$ for Sigfox Antwerp (spatially sparse 84-BS signatures)
and $5\to3$ for LoRaWAN and Indoor (already compact telemetry).

\subsection{Stage 2: EKF-Based Digital Twin}
\label{sec:dt}

Unlike static classifiers that process each observation in isolation,
the DT (DT-EKF) maintains a \textit{dynamic} estimate of
$\mathbf{s}_i(t)=[{\rm RSSI}_i(t),\,n^{\rm BS}_i(t)]^\top$
(in linear mW domain) through a recursive EKF~\cite{ref_ekf},
providing three capabilities unavailable in static methods:
(i)~adaptive noise weighting via the time-varying Kalman gain,
(ii)~statistically principled anomaly scoring via the NIS chi-squared test,
and (iii)~data-adaptive quality threshold via Youden's J.

The proposed nonlinear state transition for the RSSI component is:
\begin{equation}
s^{(1)}_i(t+1) = s^{(1)}_i(t)
  + \alpha\tanh\!\Bigl(\tfrac{\mu_i - s^{(1)}_i(t)}{\sigma_i}\Bigr)
  + w^{(1)}_t,
\label{eq:ss1}
\end{equation}
where $\mu_i$, $\sigma_i$ are the dataset RSSI mean and standard deviation
and $\alpha{=}0.3$ is the mean-reversion strength. This
Ornstein-Uhlenbeck-like tanh term captures mean-reverting channel dynamics
with saturation at the signal boundary. The base-station count
$s^{(2)}_i$ follows a standard random walk. The observation model is
$\mathbf{o}_i(t)=\mathbf{s}_i(t)+\mathbf{v}(t)$ (identity mapping),
where $\mathbf{o}_i(t)$ is the noisy channel measurement (distinct from
the VAE code $\hat{\mathbf{z}}_i(t)$). Because the transition Jacobian
$\mathbf{F}_t$ changes at every step due to the tanh nonlinearity, full
EKF linearization~\cite{ref_ekf} is required rather than a standard Kalman Filter (KF).

\textit{Anomaly detection.}
The NIS statistic is:
\begin{equation}
\xi_t = \mathbf{y}_t^\top \mathbf{S}_t^{-1} \mathbf{y}_t \sim \chi^2(2)
\label{eq:nis}
\end{equation}
under the no-fault hypothesis, where $\mathbf{y}_t$ is the EKF innovation
and $\mathbf{S}_t$ its covariance~\cite{ref_ekf}. Step $t$ is flagged
anomalous if $\xi_t > \chi^2_{0.95}(2) \approx 5.99$, controlling the
false-alarm rate at 5\%.

\textit{Quality classification.}
The EKF-estimated power $\hat{s}^{(1)}_{t|t}$ (mW) is converted to dBm
as $\widehat{\rm RSSI}_{t}=10\log_{10}(\hat{s}^{(1)}_{t|t})$.
Because EKF smoothing shifts the filtered RSSI distribution relative to
the raw signal, the classification threshold $\tau^*$ is derived on the
EKF output via Youden's J, $\tau^*=\arg\max_\tau[\mathrm{TPR}(\tau)-\mathrm{FPR}(\tau)]$,
evaluated against the median-derived pseudo-labels from
Section~\ref{sec:model}. This two-step strategy (median threshold for
reference labels, Youden's J for EKF classifier) removes systematic
smoothing bias without manual calibration.

\subsection{Stage 3: MADRL with Graph Attention Network}

\textit{State and action.}
Agent $i$'s local state at time $t$ is the VAE latent code
$\hat{\mathbf{z}}_i(t)\in\mathbb{R}^{d_z}$, where $d_z=16$ for
Sigfox Antwerp and $d_z=3$ for LoRaWAN Italy and Indoor
(see Table~\ref{tab:config}). Agents are connected in
graph $\mathcal{G}$ and each selects action $a_i(t)\in\{1,\ldots,A\}$.
The EKF quality label $\hat{q}_i(t)$ and anomaly flag $\xi_i(t)$ from
Stage~2 enter the learning signal via the reward rather than the state
vector. This keeps the state compact but means agents cannot proactively
branch on anomaly information before receiving delayed reward feedback,
a limitation identified as future work (Section~\ref{sec:experiments}).

\textit{GAT encoder.}
Agent representations are aggregated via standard GAT~\cite{ref_gat_comm},
producing attention-weighted neighborhood embeddings that allow each agent
to selectively emphasize the most relevant neighbors (e.g., strongly
interfering or heavily loaded nodes) rather than applying uniform aggregation.

\textit{DQN training.}
Each agent maintains an independent Q-network $Q_i(\cdot;\bm{\theta}_i)$
trained via the standard DQN Bellman loss with discount factor $\gamma$
and a periodically updated target network $\bm{\theta}_i^-$
to stabilize training. The per-agent reward models an adaptive
power-control trade-off: higher resource action boosts effective signal
quality when the channel is weak, but always incurs a resource cost:
\vspace{-4pt}
\begin{equation}
r_i(t) = 5\,\tilde{s}_i(t) + 2\,\bar{n}_i(t)
        - 3\,\bar{a}(t)
        + \max\!\bigl(0,\,1 - \tfrac{\sigma_a(t)}{2}\bigr),
\label{eq:reward}
\end{equation}
\vspace{-2pt}
where $\bar{s}_i(t)=\mathrm{clip}\!\bigl(({\rm RSSI}_i(t)+145)/95,0,1\bigr)\in[0,1]$
is the RSSI normalized to $[-145,-50]$~dBm, avoiding any sign ambiguity
and ensuring scale-invariant reward regardless of deployment distance,
$\tilde{s}_i(t)=\bar{s}_i(t)+\bar{a}(t)(1-\bar{s}_i(t))$ is the
effective normalized RSSI boosted by mean action
$\bar{a}(t)=(1/N)\sum_i a_i(t)/(A{-}1)\in[0,1]$
proportional to channel need $(1-\bar{s}_i)$,
$\bar{n}_i(t)=\min(n_i^{\rm BS}/20,1)$ is the normalized active-BS count,
and $\sigma_a(t)$ is the cross-agent action standard deviation.
The total reward lies in $[-3,\,8]$, requiring no additional clipping.
This design creates a data-driven cross-over: when the channel is poor
(low $\bar{s}_i$), boosting resources yields a net gain. When the
channel is good, conservative allocation is optimal.

\begin{table}[t]
\centering
\caption{Dataset Characteristics and Hyperparameters}
\label{tab:config}
\setlength{\tabcolsep}{3pt}
\begin{tabular}{lccc}
\toprule
\textbf{Parameter} & \textbf{Antwerp} & \textbf{LoRaWAN} & \textbf{Indoor}\\
\midrule
Technology           & Sigfox  & LoRaWAN & WiFi   \\
Samples $N$          & 14,377  & 190,502 & 1,100  \\
VAE input dim $d$    & 84      & 5       & 5      \\
VAE latent dim $d_z$ & 16      & 3       & 3      \\
$\sigma_Q$           & 0.5     & 0.3     & 0.2    \\
$\sigma_R$           & 0.5     & 0.4     & 0.8    \\
Agents $N$           & 5       & 5       & 5      \\
Actions $A$          & 5       & 5       & 5      \\
\bottomrule
\end{tabular}
\end{table}

\begin{algorithm}[H]
\caption{Proposed Framework: Training Procedure}
\label{alg:main}
\begin{algorithmic}[1]
\REQUIRE Dataset $\mathcal{D}$; config $\{d_z,\sigma_Q,\sigma_R,\varepsilon_0,\gamma,N_{\rm ep}\}$
\STATE // Stage 1: VAE pre-training
\STATE Normalise $\mathbf{X}\leftarrow(\mathbf{X}-\mu)/\sigma$
\FOR{epoch $e=1$ \textbf{to} $E_{\rm VAE}$}
  \STATE Minimise $-\mathcal{L}(\phi,\theta;\mathbf{X})$ via Adam
\ENDFOR
\STATE Freeze encoder; extract $\hat{\mathbf{Z}}=\bm{\mu}_\phi(\mathbf{X})$
\STATE // Stage 2: DT-EKF calibration
\STATE Run EKF on $\{\mathbf{s}_i(t)\}$ following~\cite{ref_ekf}
\STATE Compute $\tau^*$ via Youden's J
\STATE // Stage 3: MADRL-GAT training
\STATE Init $\{Q_i,\mathcal{B}_i\}_{i=1}^N$; $\varepsilon\leftarrow\varepsilon_0$
\FOR{episode $e=1$ \textbf{to} $N_{\rm ep}$}
  \FOR{each time step $t$}
    \STATE Encode graph via GAT~\cite{ref_gat_comm}
    \STATE Select $a_i$ via $\varepsilon$-greedy; observe $r_i(t)$ via (\ref{eq:reward})
    \STATE Store in $\mathcal{B}_i$; update $Q_i$ via DQN Bellman loss
  \ENDFOR
  \STATE Decay $\varepsilon$; soft-update $\bm{\theta}_i^-$
\ENDFOR
\end{algorithmic}
\end{algorithm}

% ============================================================
\section{Experimental Results}
\label{sec:experiments}
% ============================================================

\subsection{Experimental Settings}

Experiments run on a CPU-only workstation (Python 3.11, PyTorch 2.x,
scikit-learn). Table~\ref{tab:config} summarizes the per-dataset
hyperparameters. The VAE trains for 30 epochs with Adam
($\mathrm{lr}=10^{-3}$, cosine annealing). MADRL trains for 50 episodes
with $\gamma=0.99$, $\varepsilon$ decaying from 0.1 to 1.0, replay
batch size 16, and target network updated every 10 steps.


Three real-world IoT datasets spanning heterogeneous wireless
technologies are used, summarized in Table~\ref{tab:datasets}.
The \textit{Sigfox Antwerp}~\cite{ref_sigfox} dataset contains 14,377
uplink transmissions captured by 84 fixed base stations across an urban
area, producing a 84-dimensional sparse RSSI signature per packet and
representing a long-range, low-power wide-area network (LPWAN) with
strong spatial structure.
The \textit{LoRaWAN Italy}~\cite{ref_lorawan} dataset comprises 190,502
outdoor sensor readings with five features (RSSI, SNR, soil temperature,
humidity, battery voltage), covering a large-scale agricultural deployment
where channel conditions vary substantially with distance and vegetation.
The \textit{Indoor WiFi}~\cite{ref_indoor} dataset consists of 1,100
fingerprint samples collected from five access points inside a building,
where multipath propagation yields a compact but highly variable 5-dimensional
RSSI vector.
Together, the three datasets exercise the proposed framework under
contrasting dimensionality (84 vs.\ 5), scale (1,100 to 190,502 samples),
and deployment contexts (urban LPWAN, agricultural IoT, indoor WiFi).
Quality thresholds are set to the per-dataset median RSSI to ensure
approximately balanced binary class labels regardless of deployment-specific
link-budget conventions.

\begin{table}[t]
\centering
\caption{Summary of Real-World IoT Datasets}
\label{tab:datasets}
\setlength{\tabcolsep}{4pt}
\renewcommand{\arraystretch}{1.15}
\begin{tabular}{lccc}
\toprule
\textbf{Property}        & \textbf{Antwerp}           & \textbf{LoRaWAN}           & \textbf{Indoor}           \\
\midrule
Technology               & Sigfox~\cite{ref_sigfox}   & LoRaWAN~\cite{ref_lorawan} & WiFi~\cite{ref_indoor}    \\
Samples                  & 14,377                     & 190,502                    & 1,100                     \\
Input dim $d$            & 84                         & 5                          & 5                         \\
Features                 & \shortstack[c]{RSSI\\$\times$84~BSs}
                         & \shortstack[c]{RSSI, SNR,\\soil temp.,\\humid., batt.}
                         & \shortstack[c]{RSSI\\$\times$5~APs}                     \\
Quality threshold $\tau$ & $-$123.9~dBm               & $-$69.0~dBm                & $-$63.9~dBm               \\
Good-quality ratio       & 50.0\%                     & 47.9\%                     & 49.5\%                    \\
\bottomrule
\end{tabular}
\end{table}

Comparison methods.
Because the proposed system is an integrated three-stage framework, the
evaluation distinguishes two levels of comparison: \textit{module-level}
reference methods that isolate the contribution of the DT-EKF stage, and
\textit{framework-level} baselines that compare alternative resource
allocation policies against the full pipeline.

\textit{Module-level references for DT-EKF.}
To assess the quality classification capability of the DT-EKF module,
two representative unsupervised static clustering methods are used:
K-Means ($k=2$) partitions observations via hard assignment, while
Gaussian Mixture Model (GMM, $k=2$) extends this with probabilistic
soft assignments and full covariance modelling.
Setting $k=2$ directly matches the binary good/poor quality objective.
Neither method exploits temporal dynamics, so comparing them against the
DT-EKF quantifies the benefit of recursive state estimation over
snapshot-based clustering.
For anomaly detection, three classical single-task detectors are used
as module-level references:
Local Outlier Factor (LOF, $k=20$) represents density-based detection;
Isolation Forest (IF, 200 trees) represents ensemble-based detection;
and Rolling Z-Score (window $W=20$, $z_t=|x_t-\bar{x}|/\hat{\sigma}$)
represents lightweight univariate statistical detection.
These are compared against the DT-EKF NIS $\chi^2$ test to evaluate
whether multivariate probabilistic state estimation improves over
single-task classical methods.
Critically, none of these module-level references is a full-system
alternative: K-Means/GMM cannot detect anomalies, and
LOF/IF/Z-Score cannot classify quality.
The proposed DT-EKF is the only evaluated method to deliver both
capabilities simultaneously within a single recursive filter pass.

\textit{Framework-level baselines for MADRL-GAT.}
Three resource allocation policies are compared against the full
MADRL-GAT pipeline to bracket performance and isolate architectural
contributions:
Random policy (uniform random action) establishes a stochastic lower bound;
Greedy-min (always action~0, minimum power) represents a deterministic
conservative heuristic; and DQN-noGAT replaces the GAT encoder with a
plain Multilayer Perceptron (MLP), isolating the contribution of
inter-agent graph-attention coordination.

\subsection{VAE Semantic Compression}

\begin{figure*}[t]
  \centering
  \subfloat[Antwerp ($84\!\to\!16$)]{%
    \includegraphics[width=0.32\textwidth]{vae_loss_antwerp}}
  \hfill
  \subfloat[LoRaWAN ($5\!\to\!3$)]{%
    \includegraphics[width=0.32\textwidth]{vae_loss_lorawan}}
  \hfill
  \subfloat[Indoor ($5\!\to\!3$)]{%
    \includegraphics[width=0.32\textwidth]{vae_loss_indoor}}
  \caption{VAE training loss curves. Solid: reconstruction loss,
  dashed: KL divergence. Both converge within 30 epochs.}
  \label{fig:vae_loss}
\end{figure*}


Fig.~\ref{fig:vae_loss} illustrates the training dynamics of the VAE across the three datasets.
Both the reconstruction loss and the KL divergence converge rapidly, with the dominant
improvement occurring within the first few epochs. This fast convergence is a direct
consequence of the data characteristics rather than a training artifact.
For Sigfox Antwerp, the 84-dimensional RSSI signature is spatially sparse: each device
is typically received by only 2--5 out of 84 base stations, so the majority of input
dimensions are zero after normalization. The VAE quickly learns to reconstruct these
inactive dimensions to zero, which accounts for the steep initial loss drop.
For LoRaWAN and Indoor, the input dimensionality is already low (5 features), making
the $5\!\to\!3$ compression an intrinsically simple task that any small network solves
within a few gradient steps. In both cases, the KL term stabilizes rapidly because
the low-dimensional, pre-normalized data provides little information beyond what the
prior already captures, so the encoder settles at a compact posterior early in training.
The remaining epochs refine the reconstruction of the few active or informative
dimensions, yielding marginal but consistent improvement. Critically, the latent codes
produced after this rapid convergence are sufficiently informative to drive stable
MADRL-GAT learning in Stage~3, confirming that early saturation of the loss does not
imply poor representation quality.

Evaluating compression efficacy, the Sigfox Antwerp configuration attains an MSE of $0.0082$ under a severe dimension reduction factor of $5.25\times$ (from $84$ to $16$). This significant compressibility is attributed to the spatial sparsity inherent in the 84-dimensional Sigfox signatures, as transmissions are generally captured by only a limited subset of base stations. Consequently, the VAE proficiently mitigates spatial redundancy and observational noise to isolate the essential semantic features.

In contrast, the LoRaWAN and Indoor datasets, which are characterized by denser, lower-dimensional sensor telemetry, yield MSEs of $0.016$ and $0.054$, respectively, with a compression factor of $1.67\times$ ($5 \to 3$). Despite the more modest reduction ratio, the architecture robustly distills a compact and noise-tolerant semantic state. This state-space contraction is pivotal for the downstream MADRL operations (Stage 3), as it exponentially curtails the state-space dimensionality, which in turn accelerates Q-network convergence and enhances overall sample efficiency.

These results demonstrate that the VAE effectively captures the underlying semantic structure across heterogeneous IoT telemetry, providing compact yet informative latent representations that reduce state-space complexity for all downstream stages.

\subsection{Digital Twin EKF Performance}

\begin{figure*}[t]
  \centering
  \subfloat[Antwerp (RSSI-RMSE=3.48~dBm)]{%
    \includegraphics[width=0.32\textwidth]{dt_estimation_antwerp}}
  \hfill
  \subfloat[LoRaWAN (RSSI-RMSE=5.51~dBm)]{%
    \includegraphics[width=0.32\textwidth]{dt_estimation_lorawan}}
  \hfill
  \subfloat[Indoor (RSSI-RMSE=2.47~dBm)]{%
    \includegraphics[width=0.32\textwidth]{dt_estimation_indoor}}
  \caption{DT-EKF state estimation over 500 time steps.}
  \label{fig:dt}
\end{figure*}

Fig.~\ref{fig:dt} shows that the EKF effectively smooths injected faults
and tracks the underlying channel trend. RSSI-RMSE is highest for LoRaWAN
(5.51~dBm) owing to its wide dynamic range ($-20$ to $-120$~dBm).
Indoor achieves the lowest RMSE (2.47~dBm) due to spatially constrained
WiFi propagation.

Table~\ref{tab:main_results} compares all methods comprehensively.
DT-EKF outperforms K-Means and GMM on quality classification across all
datasets. The most significant improvement is on LoRaWAN, where K-Means
and GMM achieve AUC $\approx0.505$ (near-random), while DT-EKF reaches
$\text{AUC}=0.9106$ (a relative gain of +80.4\%). For Indoor, Youden's-J threshold optimization pushes accuracy to 0.934, surpassing
K-Means by +16.6 percentage points.

\begin{figure*}[t]
  \centering
  \subfloat[Antwerp]{%
    \includegraphics[width=0.48\textwidth]{roc_comparison_antwerp}}
  \hfill
  \subfloat[LoRaWAN]{%
    \includegraphics[width=0.48\textwidth]{roc_comparison_lorawan}}
  \\[0.5em]
  \centering
  \subfloat[Indoor]{%
    \includegraphics[width=0.48\textwidth]{roc_comparison_indoor}}
  \caption{ROC curves per dataset.}
  \label{fig:roc}
\end{figure*}
\begin{table*}[t]
\centering
\caption{Comprehensive Performance Comparison on Three IoT Datasets.}
\label{tab:main_results}
\setlength{\tabcolsep}{4pt}
\renewcommand{\arraystretch}{1.1}
\begin{tabular}{llccccccc}
\toprule
\multirow{2}{*}{\textbf{Dataset}} &
\multirow{2}{*}{\textbf{Method}} &
\multicolumn{2}{c}{\textbf{Quality Classification}} &
\multicolumn{4}{c}{\textbf{Anomaly Detection (AUC)}} &
\multirow{2}{*}{\textbf{VAE MSE}} \\
\cmidrule(lr){3-4}\cmidrule(lr){5-8}
 & & \textbf{Acc.} & \textbf{AUC} &
\textbf{LOF} & \textbf{IF} & \textbf{Z-Score} & \textbf{DT-EKF}$^\dagger$ & \\
\midrule
\multirow{3}{*}{Antwerp}
 & K-Means & 0.730 & 0.676 &
   \multirow{3}{*}{0.375} & \multirow{3}{*}{\textbf{0.949}} &
   \multirow{3}{*}{0.920} & \multirow{3}{*}{0.835} &
   \multirow{3}{*}{0.0082} \\
 & GMM     & 0.656 & 0.697 & & & & & \\
 & \textbf{DT-EKF} & \textbf{0.818} & \textbf{0.867} & & & & & \\
\midrule
\multirow{3}{*}{LoRaWAN}
 & K-Means & 0.505 & 0.505 &
   \multirow{3}{*}{0.500} & \multirow{3}{*}{\textbf{0.992}} &
   \multirow{3}{*}{0.919} & \multirow{3}{*}{0.796} &
   \multirow{3}{*}{0.0160} \\
 & GMM     & 0.505 & 0.516 & & & & & \\
 & \textbf{DT-EKF} & \textbf{0.846} & \textbf{0.911} & & & & & \\
\midrule
\multirow{3}{*}{Indoor}
 & K-Means & 0.767 & 0.830 &
   \multirow{3}{*}{0.307} & \multirow{3}{*}{0.898} &
   \multirow{3}{*}{0.958} & \multirow{3}{*}{\textbf{0.987}} &
   \multirow{3}{*}{0.0537} \\
 & GMM     & 0.611 & 0.666 & & & & & \\
 & \textbf{DT-EKF} & \textbf{0.934} & \textbf{0.972} & & & & & \\
\bottomrule
\end{tabular}
\end{table*}

For anomaly detection, Table~\ref{tab:main_results} reveals a
structured hierarchy among the baselines.
LOF achieves near-random AUC across all three datasets
(0.375 / 0.500 / 0.307), confirming that density-based neighbourhood
methods are poorly suited to the spike-and-drop injection pattern used
in this evaluation.
Isolation Forest (IF) is the strongest single-task baseline on Antwerp
(AUC 0.949) and LoRaWAN (0.992), with its
random sub-space isolation well-matched to these abrupt fault patterns.
Z-Score ranks second on both (0.920 / 0.919), while DT-EKF trails
(0.835 / 0.796) on these two single-dominant-channel datasets.
DT-EKF reverses the ranking on Indoor (0.987), surpassing
Z-Score (0.958) and IF (0.898), because the bivariate NIS statistic
jointly exploits complementary fault signatures in both RSSI and the
active base-station count, a structural advantage unavailable to
univariate or purely feature-based detectors.
Crucially, IF and Z-Score are dedicated single-task anomaly detectors
that provide no quality-classification capability, while DT-EKF
simultaneously produces Youden-optimal quality labels and anomaly scores
within the same recursive filter pass, without any additional model,
training data, or labels.
This dual-task capability is the defining advantage of the proposed
DT and cannot be reproduced by any single competing method.

These results establish that the DT-EKF is the strongest
multi-purpose monitoring component across the pipeline: it
consistently surpasses static clustering baselines in quality
classification by $+$28.3\%, $+$80.4\%, and $+$17.1\% AUC on
Antwerp, LoRaWAN, and Indoor, and provides competitive anomaly
detection, winning decisively on Indoor.
Among single-task anomaly detectors, IF achieves the highest AUC on
two datasets but its inability to provide quality labels or temporal
state estimates makes it a complement, not an alternative, to the
proposed unified pipeline.

\subsection{MADRL-GAT Resource Allocation}

\begin{figure*}[t]
  \centering
  \subfloat[Antwerp (MADRL-GAT Eval: 2305.05)]{%
    \includegraphics[width=0.32\textwidth]{madrl_rewards_antwerp}}
  \hfill
  \subfloat[LoRaWAN (MADRL-GAT Eval: 2519.26)]{%
    \includegraphics[width=0.32\textwidth]{madrl_rewards_lorawan}}
  \hfill
  \subfloat[Indoor (MADRL-GAT Eval: 3088.95)]{%
    \includegraphics[width=0.32\textwidth]{madrl_rewards_indoor}}
  \caption{MADRL-GAT training reward.}
  \label{fig:madrl}
\end{figure*}

Fig.~\ref{fig:madrl} shows monotonically increasing rewards across all
datasets, confirming stable policy convergence within 50 training
episodes. Table~\ref{tab:madrl_results} compares the evaluation rewards
of all resource-allocation methods.

\begin{table}[t]
\centering
\caption{Evaluation reward (mean over 20 episodes).
$\dagger$~Reward ceiling reached by all deterministic policies.}
\label{tab:madrl_results}
\renewcommand{\arraystretch}{1.1}
\resizebox{\columnwidth}{!}{%
\begin{tabular}{lcccc}
\toprule
\textbf{Dataset} & \textbf{Random} & \textbf{Greedy-min} &
\textbf{DQN-noGAT} & \textbf{MADRL-GAT (Proposed)} \\
\midrule
Antwerp  & 1813.9 & 1916.0 & 2119.0 & \textbf{2305.1} \\
LoRaWAN  & 1727.1 & 2519.3$^\dagger$ & 2519.3$^\dagger$ & \textbf{2519.3}$^\dagger$ \\
Indoor   & 2236.6 & 3088.9$^\dagger$ & 3088.9$^\dagger$ & \textbf{3088.9}$^\dagger$ \\
\bottomrule
\end{tabular}}
\end{table}

On Antwerp, MADRL-GAT achieves a reward of 2305.1, outperforming
DQN-noGAT (2119.0, $+$8.8\%), Greedy-min (1916.0, $+$20.3\%), and
Random (1813.9, $+$27.1\%), demonstrating the clear benefit of
graph-attention cooperative learning on the high-dimensional
84-feature Sigfox dataset.
On LoRaWAN and Indoor, all three deterministic policies (Greedy-min,
DQN-noGAT, MADRL-GAT) converge to the same maximum-feasible reward
(2519.3 and 3088.9 respectively), as the compact 5-feature input and
fully-connected graph leave little room for the GAT encoder to
differentiate agent strategies. Random is the only policy that
fails to saturate, confirming that any structured policy suffices
but structure is required.
The compact 3-dimensional latent state for LoRaWAN and Indoor
(vs.\ raw 5-dimensional features) accelerates Q-network convergence,
a direct benefit of Stage~1 VAE compression.

These observations confirm that the integration of VAE-compressed
semantic states with quality-aware reward shaping enables stable and
efficient policy learning, regardless of the underlying wireless
technology and input feature dimensionality.

\subsection{Discussion}

DT-EKF is the only approach in the evaluation that simultaneously
performs quality classification and anomaly detection: K-Means and GMM
provide quality labels only, while LOF, IF, and Z-Score provide anomaly
scores only.
DT-EKF achieves quality AUC improvements of $+28.3\%$, $+80.4\%$,
and $+17.1\%$ over K-Means on Antwerp, LoRaWAN, and Indoor respectively,
confirming the value of temporal state estimation over static clustering.
For anomaly detection, IF is the strongest single-task baseline on
Antwerp (0.949) and LoRaWAN (0.992), while DT-EKF is the strongest
overall on Indoor (0.987), where its bivariate NIS exploits
complementary fault signatures in both channels simultaneously.
LOF is near-random on all datasets (0.375 / 0.500 / 0.307), showing
that density-based detection is ill-suited to abrupt injection faults.
This trade-off must be weighed against a fundamental capability
asymmetry: IF and Z-Score are dedicated anomaly detectors with no
quality-classification ability, whereas DT-EKF delivers both outputs
in a single recursive filter pass with no additional model, training
data, or computational overhead which is a dual-task efficiency that
no single-purpose detector can match.

\textit{Computational overhead.}
The framework's online inference involves three lightweight stages:
a VAE encoder forward pass ($\mathcal{O}(d\,d_z)$), an EKF predict-update
cycle ($\mathcal{O}(m^2)$ for the $m{=}2$ channel state), and a GAT
Q-network pass ($\mathcal{O}(N^2 d_z)$ per head).
All operations run on CPU in under 1~ms per time step in our experiments.
Formal latency profiling on constrained edge hardware (e.g., ARM
Cortex-M or Raspberry Pi) remains a priority for future hardware-in-the-loop
validation to confirm real-time suitability on the most resource-limited IoT nodes.

\textit{Concept drift.}
The VAE encoder is trained offline on a fixed data snapshot, making its
latent codes sensitive to distribution shift caused by seasonal variation,
hardware aging, or network reconfiguration.
To mitigate this, a lightweight periodic fine-tuning strategy is proposed:
during scheduled low-traffic windows, a small batch of recent unlabeled
measurements updates only the encoder parameters via the reconstruction
loss, leaving the frozen DT and MADRL stages intact.
This continual adaptation keeps the semantic feature space aligned with
the evolving channel distribution while preserving downstream pipeline
stability, and is planned as part of a deployed edge-IoT prototype.

% ============================================================
\section{Conclusion}
\label{sec:conclusion}
% ============================================================

A unified three-stage framework for intelligent IoT resource management
has been presented, coupling VAE semantic compression, DT-EKF
joint monitoring, and MADRL-GAT cooperative resource allocation into a
sequential pipeline. VAE-compressed latent codes reduce the MADRL state
space, and EKF quality labels and anomaly flags are fed directly into
the per-agent reward, establishing a concrete information path from
sensing to resource allocation. Unlike prior approaches that address
semantic compression, quality classification, and anomaly detection in
isolation, the proposed pipeline delivers all three capabilities in a
single system validated across three wireless technologies.
The DT-EKF substantially outperforms static unsupervised
baselines in quality classification across all evaluated IoT technologies
while simultaneously detecting anomalies within the same filter pass,
a dual capability unavailable in any single competing method.
For anomaly detection, DT-EKF achieves a decisive advantage on Indoor
(AUC 0.987 vs.\ IF 0.898, Z-Score 0.958). On Antwerp and LoRaWAN,
Isolation Forest achieves higher single-channel AUC (0.949 / 0.992)
but provides no quality-classification capability, whereas DT-EKF
delivers both outputs simultaneously without additional computational cost. VAE-compressed semantic states
reduce the effective state space, enabling stable policy convergence
within 50 training episodes across all three datasets. Cross-technology
generalizability (Sigfox, LoRaWAN, WiFi) is confirmed, extending beyond
the single-domain evaluations in prior work. Future work will explore
Unscented KF for stronger nonlinear channel modeling, ablation studies
to quantify the contribution of each stage, Proximal Policy Optimization (PPO)-based agents for
improved sample efficiency, and federated deployment for
privacy-preserving distributed IoT management.

% ============================================================
\begin{thebibliography}{99}

\bibitem{ref_iot_overview}
A. Al-Fuqaha, M. Guizani, M. Mohammadi, M. Aledhari, and M. Ayyash,
``Internet of Things: A survey on enabling technologies, protocols, and
applications,''
\textit{IEEE Commun. Surv. Tutor.}, vol.~17, no.~4, pp.~2347--2376, 2015.

\bibitem{ref_drl_comm2}
Y. Xiao, Y. Song, and J. Liu,
``Multi-agent deep reinforcement learning based resource allocation for ultra-reliable low-latency Internet of controllable things,''
\textit{IEEE Trans. Wireless Commun.}, vol.~22, no.~8, pp.~5414--5430, 2023.

\bibitem{ref_youden}
W.~J. Youden,
``Index for rating diagnostic tests,''
\textit{Cancer}, vol.~3, no.~1, pp.~32--35, 1950.

\bibitem{ref_vae_wireless}
C. Liu, C. Guo, Y. Yang, and N. Jiang,
``Adaptable semantic compression and resource allocation for
task-oriented communications,''
\textit{IEEE Trans. Cogn. Commun. Netw.}, vol.~10, no.~3,
pp.~769--782, 2024.

\bibitem{ref_dt_iot}
R. Minerva, G. M. Lee, and N. Crespi,
``Digital twin in the IoT context: A survey on technical features,
scenarios, and architectural models,''
\textit{Proc. IEEE}, vol.~108, no.~10, pp.~1785--1824, 2020.

\bibitem{ref_drl_iot}
X. Xiong, K. Zheng, L. Lei, and L. Hou,
``Resource allocation based on deep reinforcement learning in IoT edge
computing,''
\textit{IEEE J. Sel. Areas Commun.}, vol.~38, no.~6,
pp.~1133--1146, 2020.

\bibitem{ref_dt_resource}
L. Tang, Y. Du, Q. Liu, J. Li, S. Li, and Q. Chen,
``Digital-twin-assisted resource allocation for network slicing in
Industry 4.0 and beyond using distributed deep reinforcement learning,''
\textit{IEEE Internet Things J.}, vol.~10, no.~19,
pp.~16989--17006, 2023.

\bibitem{ref_semadt}
B. Hazarika, O. A. Dobre, and T. Q. Duong,
``Digital twin and semantic-aware multi-agent RL for maritime search
and rescue operations,''
in \textit{Proc. IEEE PIMRC}, 2025.

\bibitem{ref_anomaly_survey}
A. De Benedictis, F. Flammini, N. Mazzocca, A. Somma, and F. Vitale,
``Digital twins for anomaly detection in the industrial Internet of
Things: Conceptual architecture and proof-of-concept,''
\textit{IEEE Trans. Ind. Inform.}, vol.~19, no.~12,
pp.~11553--11563, 2023.

\bibitem{ref_gat_comm}
P. Velickovic, G. Cucurull, A. Casanova, A. Romero, P. Li\`{o},
and Y. Bengio,
``Graph attention networks,''
in \textit{Proc. ICLR}, 2018.

\bibitem{ref_vae}
D. P. Kingma and M. Welling,
``Auto-encoding variational Bayes,''
in \textit{Proc. ICLR}, 2014.

\bibitem{ref_ekf}
D. Simon,
``Optimal State Estimation: Kalman, $H_\infty$, and Nonlinear Approaches,''
John Wiley \& Sons, 2006.

\bibitem{ref_sigfox}
R. Calvo-Palomino et al.,
``Incremental learning meets transfer learning: Application to
multi-site Sigfox LPWAN deployments,''
in \textit{Proc. ACM IoTDI}, 2021.

\bibitem{ref_lorawan}
F. Cuomo et al.,
``EXPLoRa: Extending the performance of LoRa by suitable spreading
factor allocations,''
in \textit{Proc. IEEE WiMob}, 2017.

\bibitem{ref_indoor}
J. Torres-Sospedra et al.,
``UJIIndoorLoc: A new multi-building and multi-floor database for WLAN
fingerprint-based indoor localization,''
in \textit{Proc. IEEE IPIN}, 2014.

\bibitem{ref_wali2025}
S. Wali, M.~I. Khan, and M. Imran,
``Semantic-aware reinforcement and ensemble learning for signal management
and anomaly detection in IoT systems,''
\textit{Sci. Rep.}, vol.~15, no.~1, p.~42594, 2025.

\end{thebibliography}

\end{document}
